% Canonical intrinsic operation: no independent source order is inserted.
\section{The Gamma orders produced by the existing exterior and sum maps}
\label{sec:IGO}

The source order in this section is one of the two orders already present
in the construction. The determinant size is the actual exterior degree.
Every relation coefficient remains in the resulting density.

\subsection{Original coordinates and the exact source measure}
Retain
\[
 k=4l+1,\quad l,m\geq1,\quad e=1+k(m-1),\quad
 q=e(k+1)^2,\quad c=k/2,\quad
 0<\delta<1/2,\quad\gamma>2,
\]
\[
 \chi(S)=\prod_{a,d=0}^{k}
 [S-c-(2a-k)\delta-i(2d-k)\gamma]^e,
 \qquad s\in\{1,k\},\quad b=s/2,\quad M_s=(2\pi)^{s/2}.
 \tag{IGO1}
\]
Here $s=1$ is the fixed source $\sigma$ and $s=k$ is the original
tensor Gamma source. The original line throughout is $S=c+iy$.
The change of letter from the second root index to $d$ leaves every
root and its order unchanged.

For each real $v>0$ define a density on the real coordinate by
\[
 \rho_v(y)=\frac{2^v}{4\pi\Gamma(v)}
       \left|\Gamma(v/2+iy/2)\right|^2,
 \qquad dm_{0,s}(y)=M_s\rho_b(y)\,dy.
 \tag{IGO2}
\]
Thus the full mass $M_s$ is retained. In particular
$dm_{0,1}=|\Gamma(1/4+iy/2)|^2dy/(2\pi)=d\sigma$.
For the other existing source, the recurrence of the Gamma function
at $k/4=l+1/4$ gives
\[
 dm_{0,k}(y)=\beta_k\prod_{j=0}^{l-1}
           [y^2+(2j+\tfrac12)^2]d\sigma(y),\qquad
 \beta_k=\frac{(2\pi)^{k/2}}{\sqrt2\Gamma(k/2)}.
\]
Indeed the squared Gamma recurrence contributes $4^{-l}$ times
the displayed product. The coefficient in IGO2, multiplied by
$4^{-l}$ and divided by the coefficient $1/(2\pi)$ of $\sigma$,
is $M_k2^{k/2-1-2l}/\Gamma(k/2)=\beta_k$.
This proves agreement with the original tensor source including
its full polynomial factor and mass.
For real $|t|<\pi/2$ and real $\xi$, the exact transforms are
\[
 \int_{\mathbb R}e^{ty}\rho_v(y)\,dy=(\cos t)^{-v},\qquad
 \int_{\mathbb R}e^{i\xi y}\rho_v(y)\,dy=(\cosh\xi)^{-v},
 \qquad \phi_s(t):=\int e^{ty}dm_{0,s}=M_s(\cos t)^{-b}.
 \tag{IGO3}
\]

Here is a direct justification of the source identities and the
integrability used below. The beta integral with the substitution
$u=\exp(2x)$ gives
\[
 \frac{\Gamma(v/2+iy/2)\Gamma(v/2-iy/2)}{\Gamma(v)}
  =2\int_{\mathbb R} e^{iyx}(2\cosh x)^{-v}\,dx.
\]
Consequently $\rho_v(y)=(2\pi)^{-1}
\int e^{iyx}(\cosh x)^{-v}dx$. Fourier inversion applies: the
function and its first two derivatives are integrable, and two
integrations by parts make its transform integrable. This proves
the characteristic identity and mass one in IGO3, with the displayed
Fourier convention. More strongly, for every $0<\theta<\pi/2$, move
the $x$ contour to $x+i\theta\operatorname{sgn}(y)$. The analytic
branch of $(\cosh x)^{-v}$ is fixed by its positive values on the
real axis; it exists in this strip because $\cosh x$ has no zero
there. Its modulus decays as $\exp(-v|\operatorname{Re}x|)$ on
each closed smaller strip. The vertical sides of the rectangle
therefore tend to zero and
$\rho_v(y)\leq C_{v,\theta}\exp(-\theta|y|)$.
This proves absolute convergence of every polynomial moment with
an exponential factor $e^{ty}$ when $|t|<\pi/2$. Holomorphic
continuation of the characteristic identity through this strip
then gives the first identity in IGO3. Strict positivity follows
from the zero-free Euler Gamma product in $\operatorname{Re}z>0$:
\[
 \frac1{\Gamma(z)}=z\exp(\gamma_{\!E}z)
       \prod_{n=1}^{\infty}(1+z/n)\exp(-z/n).
\]
Indeed the logarithms of the factors after their linear terms are
removed converge locally uniformly, and none of the factors vanishes
there. This product follows by taking the beta-integral limit
$\Gamma(z)=\lim_{n\to\infty}n!n^z/[z(z+1)\cdots(z+n)]$;
the linear sum is absorbed by
$\gamma_{\!E}=\lim_n(\sum_{j=1}^n1/j-\log n)$.

The natural line of the source of order $s$ is $U=b+iy$.
Its presentation on the original line is the exact translation
\[
 U=S+(b-c),\qquad
 F(S)\longmapsto F(U-(b-c)).
 \tag{IGO4}
\]
This is an algebra isomorphism, with inverse given by the opposite
translation, and its coefficient matrix on each degree filtration
is triangular with diagonal one. In particular, using $s=1$ in
IGO2 retains the original line $c+i\mathbb R$ through IGO4.
The original tensor identification is exact in the real coordinate:
\[
 dm_{0,k}=(dm_{0,1})^{*k}.
\]
Indeed the characteristic function of the convolution on the right
is $[M_1(\cosh\xi)^{-1/2}]^k
=M_k(\cosh\xi)^{-k/2}$, which is the characteristic function
of the left side by IGO3. The finite-measure Fourier uniqueness
argument given after IGO17 proves equality. The natural individual
coordinates are $1/2+iy_j$, whose sum is $c+i\sum_jy_j$.
If the individual sources are first written on the fixed original
line as $S_j=c+iy_j$, the same sum is obtained by the exact map
$\sum_jS_j\mapsto\sum_jS_j+k(1/2-c)$, applying IGO4 to all
$k$ factors. This retains the full tensor mass $M_1^k=M_k$
and its original centre $c$.

\subsection{The actual relation exterior vector and its sum measure}
The relation multiplier is the full polynomial
\[
 W_\chi(y)=|\chi(c+iy)|^2
 =\prod_{a,d=0}^{k}
 \left([(2a-k)\delta]^2+[y-(2d-k)\gamma]^2\right)^e
 =\sum_{j=0}^{2q}w_jy^j.
 \tag{IGO5}
\]
All its original parameters occur in the displayed product.
It is real, even, strictly positive on $\mathbb R$, and monic of
degree $2q$. Evenness follows by the permutation $d\mapsto k-d$.
Strict positivity follows because $k$ is odd, so
$(2a-k)\delta\ne0$ for every $a$. Thus $w_{2q}=1$ and
$w_j=0$ for odd $j$; no other coefficient is removed.

For $r\geq1$ put
\[
 \Delta_r(y)=\prod_{1\leq i<j\leq r}(y_j-y_i),\qquad
 d\lambda_{\chi,s,r}(y_1,\ldots,y_r)
 =\frac1{r!}\Delta_r(y)^2
       \prod_{i=1}^r W_\chi(y_i)\,dm_{0,s}(y_i),
 \qquad \Sigma_r(y)=\sum_{i=1}^r y_i,
 \quad \nu_{\chi,s,r}=(\Sigma_r)_*\lambda_{\chi,s,r}.
 \tag{IGO6}
\]
The notation with subscript $\varnothing$ instead of $\chi$ means
that each factor $W_\chi$ is replaced by the literal factor $1$.
For $r=0$, the product space is the one-point space with mass one,
the sum is $0$, and both pushforward measures are $\delta_0$.

This is precisely the existing exterior measure. To check all
phases and factorials, write $S_i=c+iy_i$ and define the exterior
vector in the product source Hilbert space by
\[
 \psi_{\chi,r}(y)=\frac1{\sqrt{r!}}
           \det[\chi(S_i)S_i^{j-1}]_{i,j=1}^r
   =\frac{i^{r(r-1)/2}}{\sqrt{r!}}
           \Delta_r(y)\prod_{i=1}^r\chi(c+iy_i).
 \tag{IGO7}
\]
Its squared modulus times $\prod dm_{0,s}(y_i)$ is exactly IGO6.
The phase in IGO7 is retained in the exterior map below.
Expanding both determinants and relabelling integration variables
shows that $\|\psi_{\chi,r}\|^2$ equals the Gram determinant of
$\chi,\chi S,\ldots,\chi S^{r-1}$ in the original source.
The $r!$ identical permutation contributions cancel the displayed
$1/r!$. This also proves the following full Laplace identity:
\[
 g_{\chi,s}(t)=\int e^{ty}W_\chi(y)\,dm_{0,s}(y)
       =W_\chi(\partial_t)\phi_s(t),\qquad
 H_{\chi,s,r}(t):=\int e^{tY}\,d\nu_{\chi,s,r}(Y)
       =\det[g_{\chi,s}^{(i+j)}(t)]_{i,j=0}^{r-1}.
 \tag{IGO8}
\]
All expansions and differentiations here are absolutely convergent
by the exponential estimate after IGO3, since the remaining
factors are fixed polynomials. For real $|t|<\pi/2$ the matrix
on the right is the Gram matrix of $1,y,\ldots,y^{r-1}$ for the
strictly positive measure $e^{ty}W_\chi(y)dm_{0,s}(y)$.
It is positive definite: a nonzero polynomial is nonzero on a real
open interval. Thus $H_{\chi,s,r}(t)>0$.

\subsection{The exact differential polynomial and its degree}
Put $z=\tan t$ and define polynomials recursively by
\[
 P_0(z)=1,\qquad
 P_{j+1}(z)=bzP_j(z)+(1+z^2)P_j'(z),\qquad
 Q_0(z)=\sum_{j=0}^{2q}w_jP_j(z),\qquad
 Q_{j+1}(z)=bzQ_j(z)+(1+z^2)Q_j'(z).
 \tag{IGO9}
\]
The product rule proves
$\phi_s^{(j)}(t)=\phi_s(t)P_j(\tan t)$ and
$g_{\chi,s}^{(j)}(t)=\phi_s(t)Q_j(\tan t)$.
These are exact polynomial recursions in all the $w_j$, with no
limit or discarded term.

For polynomials $F_0,\ldots,F_{r-1}$ write
$\operatorname{Wr}_z(F_0,\ldots,F_{r-1})
=\det[\partial_z^iF_j]_{i,j=0}^{r-1}$. Then
\[
 B=br+r(r-1),\qquad
 D_{\chi,s,r}(z)=\operatorname{Wr}_z(Q_0,\ldots,Q_{r-1}),
 \qquad
 H_{\chi,s,r}(t)=M_s^r(\cos t)^{-B}D_{\chi,s,r}(\tan t).
 \tag{IGO10}
\]
To prove this without losing a determinant factor, the matrix in
IGO8 is the $t$-Wronskian of $g,g',\ldots,g^{(r-1)}$.
In a Wronskian, multiplying every column function by a common
function $h$ multiplies the determinant by $h^r$:
the product-rule matrix on its rows is lower triangular with
diagonal $h$. Changing the independent variable from $t$ to $z$
multiplies its rows by a lower triangular matrix with diagonal
$1,z',\ldots,(z')^{r-1}$. Its determinant is
$(z')^{r(r-1)/2}$. Here $h=\phi_s$ and
$z'=1+z^2=(\cos t)^{-2}$. Multiplication of these two factors
gives exactly $M_s^r(\cos t)^{-br-r(r-1)}$ in IGO10.

The polynomial $P_j$ has parity $j$, degree $j$, and leading
coefficient $(b)_j=\prod_{h=0}^{j-1}(b+h)$; this follows directly
from IGO9. The polynomial $Q_j$ has parity $j$, degree $2q+j$,
and leading coefficient $(b)_{2q}(b+2q)_j$. Therefore
\[
 D_{\chi,s,r}(z)=\sum_{j=0}^{qr}d_jz^{2j},\qquad
 d_{qr}=((b)_{2q})^r
             \prod_{j=0}^{r-1}j!(b+2q)_j>0,
 \qquad D_{\chi,s,r}(z)>0\quad(z\in\mathbb R).
 \tag{IGO11}
\]
For the degree and leading coefficient, factor the leading
coefficient of each $Q_j$ from the Wronskian. If its degrees are
$n_j=2q+j$, the coefficient of its largest possible power of $z$
is $\det[(n_j)_{\underline i}]_{i,j=0}^{r-1}$.
The falling factorial in row $i$ is a monic polynomial of degree
$i$ in $n_j$; subtracting lower rows therefore changes it to the
ordinary Vandermonde determinant. Its value is
$\prod_{i<j}(n_j-n_i)=\prod_{j=0}^{r-1}j!$.
The power is $\sum_j(2q+j)-\sum_i i=2qr$, so the coefficient
cannot cancel. The parity of entry $(i,j)$ is $i+j$ modulo two;
each determinant term has even total parity
$\sum_i i+\sum_j j=r(r-1)$. Finally positivity of $D$ for
every real $z$ follows from the positive Gram in IGO8 by the
bijection $z=\tan t$ between $(-\pi/2,\pi/2)$ and $\mathbb R$.

\subsection{The Gamma order actually produced without relation factors}
For the literal exterior of the source alone, use $Q_0=1$ and
$Q_j=P_j$ in the same Wronskian calculation. Its degrees are
$j$, so its Wronskian is the positive constant
\[
 C_{b,r}=\prod_{j=0}^{r-1}j!(b)_j,\qquad
 H_{\varnothing,s,r}(t)=M_s^r C_{b,r}(\cos t)^{-B},
 \qquad
 d\nu_{\varnothing,s,r}(Y)=M_s^r C_{b,r}\rho_B(Y)\,dY.
 \tag{IGO12}
\]
The transform identity and Fourier uniqueness prove the density
formula, including its whole mass $M_s^r C_{b,r}$. The source
order in this sum law is forced by the actual exterior degree:
\[
 2B=sr+2r(r-1),\qquad
 d\nu_{\varnothing,s,r}
    =\frac{M_s^r C_{b,r}}{(2\pi)^B}\,dm_{0,2B}.
 \tag{IGO13}
\]
Here $dm_{0,a}(Y):=(2\pi)^{a/2}\rho_{a/2}(Y)dY$ names
the Gamma measure at the produced positive order $a=2B$;
the same notation below is used for the produced orders $2B+4h$.
The input order remains the original $s\in\{1,k\}$.
In the second expression $dm_{0,2B}$ is written in the real
coordinate $Y$. Its natural complex line is $U=B+iY$, whereas
the original sum coordinate is $T=rc+iY$. The exact translation
of these lines and their polynomial algebras is
\[
 U=T+(B-rc),\qquad
 F(T)\longmapsto F(U-(B-rc)).
 \tag{IGO14}
\]
It has inverse translation by $rc-B$, preserves every degree,
and has coefficient determinant one on every finite degree
filtration. If IGO4 is first applied to all $r$ factors, their
sum is $rb+iY$. The further displacement in IGO14 is exactly
$r(r-1)$, since $B-rb=r(r-1)$. Thus both the original sum centre
and the naturally produced Gamma centre have explicit maps.
Equations IGO12--IGO14 concern $r\geq1$; the unit case $r=0$
remains the point mass specified after IGO6.

\subsection{All relation terms: finite signed Gamma decomposition}
For the actual polynomial $W_\chi$, define every coefficient
\[
 a_h=\sum_{j=h}^{qr}(-1)^{j+h}\binom{j}{h}d_j,
 \qquad 0\leq h\leq qr.
 \tag{IGO15}
\]
Analytic continuation of IGO10 to $t=i\xi$, or the characteristic
integral in IGO8 directly, gives
\[
 \widehat\nu_{\chi,s,r}(\xi)
 =M_s^r(\cosh\xi)^{-B}D_{\chi,s,r}(i\tanh\xi)
 =M_s^r\sum_{h=0}^{qr}a_h(\cosh\xi)^{-(B+2h)}.
 \tag{IGO16}
\]
Indeed each term $d_j(i\tanh\xi)^{2j}$ is
$d_j(-1)^j[1-(\cosh\xi)^{-2}]^j$; expanding the finite
binomial power gives IGO15 with its stated signs. Thus
\[
 d\nu_{\chi,s,r}(Y)
  =M_s^r\sum_{h=0}^{qr}a_h\rho_{B+2h}(Y)\,dY
  =\sum_{h=0}^{qr}\frac{M_s^r a_h}{(2\pi)^{B+2h}}
                   \,dm_{0,\,2B+4h}(Y).
 \tag{IGO17}
\]
Fourier uniqueness applies to these finite signed measures. For
completeness, a finite signed measure with zero Fourier transform
has zero convolution with every Gaussian, by Fourier inversion
and Fubini. Letting the Gaussian widths decrease to zero and
testing against compactly supported continuous functions gives
zero measure. This proves the asserted equality of measures,
without a sign assumption on any $a_h$.

For each term in the second expression of IGO17 the original
sum line is still $T=rc+iY$. Its natural source line is
$U_h=B+2h+iY$, and its pullback is the explicit translation
$U_h=T+(B+2h-rc)$, with the polynomial map and inverse as in
IGO14. Thus IGO17 adds measures on the same original line; the
different natural centres have not been identified implicitly.

The top coefficient satisfies $a_{qr}=d_{qr}>0$. The largest
order forced by the full relation operation is consequently
\[
 2B+4qr=sr+2r(r-1)+4qr.
 \tag{IGO18}
\]
It equals $sq+6q^2-2q$ at $r=q$, and
$s(q+1)+6q(q+1)$ at $r=q+1$. These orders result from IGO6
and the leading coefficient of its original polynomial.
All the lower signed constituents in IGO17 remain part of the
same exact measure.

\subsection{A positive polynomial density with all masses retained}
The recurrence $\Gamma(z+1)=z\Gamma(z)$ gives, with no residual
power of two,
\[
 \frac{\rho_{B+2h}(Y)}{\rho_B(Y)}
 =\frac{\Gamma(B)}{\Gamma(B+2h)}
       \prod_{j=0}^{h-1}\left[Y^2+(B+2j)^2\right].
 \tag{IGO19}
\]
In detail, the ratio of the constants in IGO2 is
$2^{2h}\Gamma(B)/\Gamma(B+2h)$; the squared Gamma ratio is
$\prod_{j<h}[(B/2+j)^2+Y^2/4]$. Each of its $h$ factors
contributes $1/4$, cancelling $2^{2h}$ exactly. Define
\[
 R_{\chi,s,r}(X)=\sum_{h=0}^{qr}
   a_h\frac{\Gamma(B)}{\Gamma(B+2h)}
      \prod_{j=0}^{h-1}[X+(B+2j)^2].
 \tag{IGO20}
\]
Then the exact positive measure, in both conventions for the
underlying Gamma mass, is
\[
 \begin{split}
 d\nu_{\chi,s,r}(Y)
    &=M_s^r\rho_B(Y)R_{\chi,s,r}(Y^2)\,dY\\
    &=\frac{M_s^r}{(2\pi)^B}
           R_{\chi,s,r}(Y^2)\,dm_{0,2B}(Y),\\
 \deg R_{\chi,s,r}&=qr,\qquad
 [X^{qr}]R_{\chi,s,r}
       =d_{qr}\frac{\Gamma(B)}{\Gamma(B+2qr)}>0.
 \end{split}
 \tag{IGO21}
\]
The source on the second line has the centre map IGO14. In its
complex coordinate the multiplier is exactly
$R_{\chi,s,r}(-(U-B)^2)$; on the original sum line it is
$R_{\chi,s,r}(-(T-rc)^2)$.

To prove its strict positivity at every real $Y$, as opposed to
positivity only almost everywhere from the measure identity,
write the density of IGO6 relative to $dy_1\cdots dy_r$ as
\[
 w_{\chi,s,r}(y)
  =\frac{M_s^r}{r!}\Delta_r(y)^2
                  \prod_{i=1}^r W_\chi(y_i)\rho_b(y_i).
\]
For $r\geq2$ the coordinate change
$y_j=u_j$ for $j<r$ and $y_r=Y-\sum_{j<r}u_j$ has
absolute Jacobian one. It gives the actual fibre density
\[
 p_{\chi,s,r}(Y)=\int_{\mathbb R^{r-1}}
 w_{\chi,s,r}(u_1,\ldots,u_{r-1},Y-\textstyle\sum_{j<r}u_j)
                 \,du_1\cdots du_{r-1}.
 \tag{IGO22}
\]
It is finite and continuous. On a compact set of $Y$, the
exponential estimate following IGO3 bounds the integrand by
an integrable polynomial in $|u|$ times
$\exp(-\theta\sum_{j<r}|u_j|)$, with one fixed
$0<\theta<\pi/2$. This also justifies dominated convergence.
For each $Y$, choose any $t_0\ne0$ and distinct coordinates
$y_i=Y/r+t_0(i-(r+1)/2)$. Their sum is $Y$.
The Vandermonde is nonzero there, and all the other factors
are strictly positive. They remain positive in an open set
of the $u$ variables. Hence $p_{\chi,s,r}(Y)>0$.
For $r=1$ the same conclusion follows from
$p_{\chi,s,1}(Y)=M_s W_\chi(Y)\rho_b(Y)>0$.
The density in IGO21 is continuous and equals this density
almost everywhere by IGO17. Their equality therefore holds
everywhere. In particular
\[
 R_{\chi,s,r}(X)>0\quad(X\geq0),\qquad
 \nu_{\chi,s,r}(\mathbb R)
 =M_s^rD_{\chi,s,r}(0)
 =M_s^r\sum_{h=0}^{qr}a_h
 =M_s^r\int\rho_B(Y)R_{\chi,s,r}(Y^2)\,dY.
 \tag{IGO23}
\]
The coefficient sum equals $D(0)$ by the binomial identity in
IGO16 at $\xi=0$. Thus the full original relation determinant
is present as this mass. The Gamma order in IGO13, the full
polynomial in IGO20, and the signed constituents in IGO17 are
related by the proved identities IGO19--IGO23. In particular,
no choice of a larger source parameter is needed to form
this operation, and the polynomial coefficients are required
to recover its actual mass.

\subsection{The complete sum map, its adjoint, and the retained kernel}
Let $\mathcal H_{\lambda}=L^2(\mathbb R^r,\lambda_{\chi,s,r})$
and $\mathcal H_{\nu}=L^2(\mathbb R,\nu_{\chi,s,r})$, using
the full measures IGO6. The sum pullback
\[
 U_r:\mathcal H_{\nu}\longrightarrow\mathcal H_{\lambda},
 \qquad (U_rF)(y)=F(\Sigma_r y)
 \tag{IGO24}
\]
is an isometry, since its squared norm is exactly the defining
pushforward integral. Its range is closed. We use complex
inner products linear in the second variable. For $r\geq2$
its adjoint is the exact fibre integral
\[
 (U_r^*h)(Y)=\frac{1}{p_{\chi,s,r}(Y)}
  \int_{\mathbb R^{r-1}}
    h(u,Y-\textstyle\sum u_j)
    w_{\chi,s,r}(u,Y-\textstyle\sum u_j)\,du.
 \tag{IGO25}
\]
Every factor in this formula is specified in IGO22. The
denominator is the actual positive fibre mass; the source
measure and its total mass in IGO23 are unchanged. In the
Euclidean surface measure on $\Sigma_r^{-1}(Y)$, the same
numerator is $r^{-1/2}\int_{\Sigma_r^{-1}(Y)}hw\,d\mathcal H^{r-1}$.
Indeed $|\nabla\Sigma_r|=\sqrt r$, and the parametrization
in IGO22 has surface Jacobian $\sqrt r$.

For a proof of the adjoint formula and its domain, Cauchy--Schwarz
on each fibre gives
$|\int hw\,du|^2\leq p(Y)\int|h|^2w\,du$.
After division by $p(Y)$ and integration in $Y$, this gives
$\|U_r^*h\|_{\nu}\leq\|h\|_{\lambda}$, so the formula is
well defined for every equivalence class in $\mathcal H_\lambda$.
Fubini applied first to bounded integrable functions proves
$\langle U_rF,h\rangle_\lambda=\langle F,U_r^*h\rangle_\nu$;
the bound extends it to all $L^2$ functions. Taking $h=U_rF$
in IGO25 gives $U_r^*U_rF=F$. Consequently the exact orthogonal
projection and complete kernel are
\[
 \begin{split}
 \Pi_r&=U_rU_r^*,\qquad
 \mathcal H_\lambda=U_r\mathcal H_\nu\mathbin{\oplus}\ker U_r^*,\\
 \ker U_r^*&=\left\{h\in\mathcal H_\lambda:
    \int_{\mathbb R^{r-1}}h(u,Y-\textstyle\sum u_j)
     w_{\chi,s,r}(u,Y-\textstyle\sum u_j)\,du=0
      \text{ for almost every }Y\right\},\\
 \|h\|_\lambda^2&=\|U_r^*h\|_\nu^2+
                   \|h-U_rU_r^*h\|_\lambda^2.
 \end{split}
 \tag{IGO26}
\]
This identifies both the map and every discarded direction
in a passage to the sum. For $r=1$, $\Sigma_1$ is the identity,
$U_1$ and $U_1^*$ are the identity between the identical spaces,
and the kernel is zero. For $r=0$ the two Hilbert spaces are
the one-dimensional spaces at their unit points, the maps
are the identity under that identification, and the kernel
is again zero.

For $r\geq2$ the kernel remains nonzero even in the symmetric
subspace. To see this with its exact fibre dependence, let
$A(y)=\sum_i y_i^2$ and let
\[
 \overline A(Y)=\frac{1}{p(Y)}
     \int_{\mathbb R^{r-1}}A(u,Y-\textstyle\sum u_j)
                 w(u,Y-\textstyle\sum u_j)\,du.
\]
All its moments here exist by the same exponential bound.
The symmetric function $A-U_r\overline A$ lies in
$\ker U_r^*$ by IGO25. It is nonzero: for each fixed $Y$,
$A$ is a nonconstant quadratic function on the sum fibre,
and the weight is positive on the open set of distinct
coordinates. It therefore has strictly positive fibre
variance. This also records a concrete direction that the
sum coordinate alone does not reconstruct.

Finally, for $r\geq1$, the exterior Hilbert space has a complete typed map
from the same sum space. Realize $\bigwedge^rL^2(dm_{0,s})$
as the antisymmetric subspace of the product $L^2$ space with
the determinant convention in IGO7. Multiplication by
$\psi_{\chi,r}$ gives a unitary map from the symmetric
subspace of $\mathcal H_\lambda$ to that antisymmetric
subspace: its norm identity is IGO7, and its inverse is
division by $\psi_{\chi,r}$ away from the diagonal. The
diagonal has product source measure zero, and $\chi$ has
no zero on the original line, so this inverse exists almost
everywhere. Division takes an antisymmetric function to a
symmetric function and the norm identity proves its $L^2$
membership. Thus
\[
 \mathcal J_r:\mathcal H_\nu\longrightarrow
                  \bigwedge^r L^2(dm_{0,s}),\qquad
 \mathcal J_rF=\psi_{\chi,r}\,F(\Sigma_r y)
 \tag{IGO27}
\]
is an isometry with the original determinant phase retained.
For an antisymmetric vector $v$, its exact adjoint is
\[
 (\mathcal J_r^*v)(Y)=\frac1{p(Y)}
 \int_{\mathbb R^{r-1}}
   \overline{\psi_{\chi,r}(u,Y-\textstyle\sum u_j)}
   v(u,Y-\textstyle\sum u_j)
   \prod_{i=1}^r[M_s\rho_b(y_i)]\,du,
 \quad y=(u,Y-\textstyle\sum u_j).
 \tag{IGO28}
\]
Its kernel is precisely the vectors for which the numerator
in IGO28 vanishes for almost every $Y$; equivalently it is
$\psi_{\chi,r}$ times the symmetric part of the kernel in
IGO26. The proof is multiplication of IGO25 by the exact
identity $w=|\psi_{\chi,r}|^2\prod_i[M_s\rho_b(y_i)]$.
It also gives the full orthogonal decomposition with
projection $\mathcal J_r\mathcal J_r^*$. The nonzero
symmetric kernel vector just constructed supplies a nonzero
exterior kernel vector through this unitary multiplication.
At $r=0$, the empty determinant is $1$ and the zeroth exterior
space is $\mathbb C$; the exterior map and its adjoint are the
identity between the unit spaces, as for IGO24--IGO26.

For a bounded measurable function $f$, multiplication by
$f(Y)$ on $\mathcal H_\nu$ intertwines through IGO24 and
IGO27 with multiplication by $f(\Sigma_r y)$ on their
target spaces. For the unbounded original sum coordinate,
the domain on $\mathcal H_\nu$ is
$\{F:\int |YF(Y)|^2d\nu(Y)<\infty\}$; the identical
integral on the target proves the same intertwining on
exactly its transported domain. Adding the original
constant centre $rc$ and the factor $i$ gives the
intertwining for $T=rc+iY$ as well. Together with the
translations in IGO14 and after IGO17, these formulas
specify the spaces, measures, actions, adjoints, and kernels
through which the intrinsically produced Gamma density
enters the original spectral-sum construction.
